
\documentclass[tikz,border=2pt]{standalone}
\usepackage[utf8]{inputenc}
\usepackage[T1]{fontenc}
\usepackage[french]{babel}
\usepackage{amsmath,amssymb}
\usepackage[table]{xcolor}
\usepackage{siunitx}
\usepackage[version=4]{mhchem}
\usepackage{tabularx,booktabs,array}
\usepackage{tikz}
\usetikzlibrary{arrows.meta,positioning,calc,fit,decorations.pathmorphing,patterns,shapes.geometric}
\usepackage{pgfplots}
\pgfplotsset{compat=1.18}
\definecolor{primary}{RGB}{14,116,144}
\definecolor{accent}{RGB}{16,185,129}
\definecolor{dark}{RGB}{15,23,42}
\definecolor{bglight}{RGB}{241,245,249}
\definecolor{borderlight}{RGB}{203,213,225}
\definecolor{examplepink}{RGB}{236,72,153}
\definecolor{remarkblue}{RGB}{14,165,233}
\definecolor{notepurple}{RGB}{99,102,241}
\definecolor{synthgreen}{RGB}{5,150,105}
\definecolor{synthorange}{RGB}{249,115,22}
\definecolor{synthviolet}{RGB}{79,70,229}
\definecolor{primary}{RGB}{14,116,144}
\definecolor{accent}{RGB}{16,185,129}
\definecolor{dark}{RGB}{15,23,42}
\definecolor{bglight}{RGB}{241,245,249}
\definecolor{borderlight}{RGB}{203,213,225}
\definecolor{examplepink}{RGB}{236,72,153}
\definecolor{remarkblue}{RGB}{14,165,233}
\definecolor{notepurple}{RGB}{99,102,241}
\definecolor{synthgreen}{RGB}{5,150,105}
\definecolor{synthorange}{RGB}{249,115,22}
\definecolor{synthviolet}{RGB}{79,70,229}
\definecolor{primary}{RGB}{14,116,144}
\definecolor{accent}{RGB}{16,185,129}
\definecolor{dark}{RGB}{15,23,42}
\definecolor{bglight}{RGB}{241,245,249}
\definecolor{examplepink}{RGB}{236,72,153}
\definecolor{remarkblue}{RGB}{14,165,233}
\definecolor{notepurple}{RGB}{99,102,241}
\definecolor{synthgreen}{RGB}{5,150,105}
\definecolor{synthorange}{RGB}{249,115,22}
\definecolor{synthviolet}{RGB}{79,70,229}
\begin{document}
\begin{tikzpicture}[scale=0.65, >=Stealth]
% Axe optique
\draw[->] (-2,0) -- (15,0) node[right]{\small axe optique};
% Objectif L1
\draw[thick, blue] (0,-2.5) -- (0,2.5);
\draw[blue, <->] (-0.3,2.2) -- (0.3,2.2);
\node[above, blue] at (0,2.6) {$L_1$ (objectif)};
\node[below, blue] at (0,-2.7) {$O_1$};
% Foyers objectif
\filldraw[blue] (5,0) circle (2pt) node[below]{\small $F'_1 = F_2$};
% Oculaire L2
\draw[thick, red] (10,-1.5) -- (10,1.5);
\draw[red, <->] (9.7,1.2) -- (10.3,1.2);
\node[above, red] at (10,1.6) {$L_2$ (oculaire)};
\node[below, red] at (10,-1.7) {$O_2$};
% Foyer image oculaire
\filldraw[red] (15,0) circle (0pt);
% Rayons entrants (parallèles inclinés vers le bas)
\draw[->, thick, orange] (-2,1.5) -- (0,1.2);
\draw[->, thick, orange] (-2,0.9) -- (0,0.6);
\draw[->, thick, orange] (-2,0) -- (0,-0.3);
% Vers image intermédiaire
\draw[thick, orange] (0,1.2) -- (5,-0.8);
\draw[thick, orange] (0,0.6) -- (5,-0.8);
\draw[thick, orange] (0,-0.3) -- (5,-0.8);
% Image intermédiaire
\draw[thick, green!60!black, ->] (5,0) -- (5,-0.8) node[right]{\small $B_1$};
\node[above right, green!60!black] at (5,0.1) {\small $A_1$};
% Rayons après oculaire (divergent de A1B1, passent par L2)
\draw[thick, purple] (5,-0.8) -- (10,-2.2);
\draw[thick, purple] (5,-0.8) -- (10,0);
% Rayons émergents (parallèles, angle plus grand)
\draw[->, thick, purple] (10,-2.2) -- (13,-3.4);
\draw[->, thick, purple] (10,0) -- (13,-1.2);
% Annotations angles
\draw[dashed, gray] (-2,0) -- (0,0);
\node at (-1.2,1) {\small $\theta$};
\draw[dashed, gray] (10,0) -- (13,0);
\node at (12.2,-0.5) {\small $\theta'$};
\end{tikzpicture}
\end{document}
